% !PRINTED 229
\weberpage{229}

%% [running head, not body text -- commented out 2026-08-07] Page furniture is
%% excluded from this edition by standing policy.
%% \S.\,73. \hfill Die Bezoutiante.

\[
\begin{array}{rcl|r}
E_{0,2} &=& - a_3\,t_2 - a_4\,t_1 & 4\,a_0\,\tau_1 \\
E_{1,2} &=& - a_3\,t_1 - a_4\,t_0 & 3\,a_1\,\tau_1 \\
E_{2,2} &=& a_1\,t_2 + a_2\,t_1 & 2\,a_2\,\tau_1 \\
E_{3,2} &=& a_0\,t_2 + a_1\,t_1 + a_2\,t_0 & a_3\,\tau_1 \\
\hline
E_{0,3} &=& - a_4\,t_2 & 4\,a_0\,\tau_0 \\
E_{1,3} &=& - a_4\,t_1 & 3\,a_1\,\tau_0 \\
E_{2,3} &=& - a_4\,t_0 & 2\,a_2\,\tau_0 \\
E_{3,3} &=& a_1\,t_2 + a_2\,t_1 + a_3\,t_0 & a_3\,\tau_0.
\end{array}
\]

Man hat diese Ausdrücke mit den rechts daneben gesetzten
Factoren zu multipliciren und zu addiren, dann noch
\[
\begin{array}{r}
-\dfrac{1}{4}\,\varphi(t)\,\varphi(\tau) \\[1ex]
= -\dfrac{1}{4}\,.(3\,a_1\,t_2 + 2\,a_2\,t_1 + a_3\,t_0)\,(3\,a_1\,\tau_2 + 2\,a_2\,\tau_1 + a_3\,\tau_0)
\end{array}
\]
hinzuzufügen und die Co\"efficienten von $\tau_h t_k$ aufzusuchen. So
ergiebt sich
\weq{14}{
\begin{array}{rl@{\qquad}rl}
B_{2,2} &= -\dfrac{1}{4}(8\,a_0\,a_2 - 3\,a_1^2), & B_{0,1} &= -\dfrac{1}{2}\;(6\,a_1\,a_4 - a_2\,a_3) \\[1.5ex]
B_{1,1} &= -4\,a_0\,a_4 - 2\,a_1\,a_3 + a_2^2, & B_{1,2} &= -\dfrac{1}{2}\;(6\,a_0\,a_3 - a_1\,a_2) \\[1.5ex]
B_{0,0} &= -\dfrac{1}{4}(8\,a_2\,a_4 - 3\,a_3^2), & B_{0,2} &= -\dfrac{1}{4}(16\,a_0\,a_4 - a_1\,a_3).
\end{array}
}
Auf dieselbe Weise kann man bei einer Form fünften Grades
verfahren, und findet für die Co\"efficienten der Bezoutiante folgende Ausdrücke
\[
\begin{array}{rcl@{\qquad}rcl}
5\,B_{0,0} &=& 4\,a_4^2 - 10\,a_3\,a_5 \\
5\,B_{1,1} &=& 6\,a_3^2 - 10\,a_2\,a_4 - 20\,a_1\,a_5 \\
5\,B_{2,2} &=& 6\,a_2^2 - 10\,a_1\,a_3 - 20\,a_0\,a_4 \\
5\,B_{3,3} &=& 4\,a_1^2 - 10\,a_0\,a_2 \\
\end{array}
\]
\[
\begin{array}{rcl}
5\,B_{0,1} &=& 3\,a_3\,a_4 - 15\,a_2\,a_5 \\
5\,B_{0,2} &=& 2\,a_2\,a_4 - 20\,a_1\,a_5 \\
5\,B_{0,3} &=& a_1\,a_4 - 25\,a_0\,a_5 \\
5\,B_{1,2} &=& 4\,a_2\,a_3 - 15\,a_1\,a_4 - 25\,a_0\,a_5 \\
5\,B_{1,3} &=& 2\,a_1\,a_3 - 20\,a_0\,a_4 \\
5\,B_{2,3} &=& 3\,a_1\,a_2 - 15\,a_0\,a_3.
\end{array}
\]
%CONT>>
